%======================================================================
\section{Conclusion}
\label{sec:conclusion}
%======================================================================

We presented the first large-scale empirical study of security vulnerabilities in LLM-generated smart contracts. Our analysis of 2,400 contracts across eight LLMs reveals that 54.7\% of compilable contracts contain at least one vulnerability, with 26.8\% containing high-severity flaws. Compilation rates range from 91.0\% (GPT-5) to 34.7\% (CodeLlama), and when controlling for contract size, vulnerability density ranges from 0.99 (GPT-5) to 4.88 (Codestral) vulns/100 LOC.

Our most striking finding is the adversarial null result: adversarial prompts do not increase vulnerability rates ($p = 0.965$, Cliff's $\delta = -0.056$), implying that vulnerabilities are \emph{intrinsic} to LLM code generation rather than prompt-dependent. We also identified three novel LLM-specific vulnerability patterns---phantom guards, inherited vulnerabilities, and context window degradation---that require LLM-aware detection tools.

Comparison with 241 verified Ethereum mainnet contracts reveals that LLM-generated vulnerability density is comparable to human-written code (2.69 vulns/100 LOC for human vs.\ 0.99--4.88 for LLM). The primary risk is not inferior code quality but the security infrastructure developers bypass when trusting LLM output.

We release our full dataset, analysis pipeline, and vulnerability taxonomy. Future work includes extending to additional LLMs and languages (Rust/Solana, Move/Aptos), runtime analysis through fuzzing, multi-turn generation sessions, and LLM-aware security tool development.
